%! Rates of Change — Unit 1, Lesson 1.3 — Group Activity (Answer Key)
\documentclass[10pt]{article}
\usepackage{precalculus-article}
\usepackage{precalculus-key}   % pulls in -boxes; do NOT also load -boxes
\newcommand{\TallMath}[1]{$\displaystyle #1\rule[-1.4em]{0pt}{3.2em}$}

\begin{document}

\pageheader{Unit 1, Lesson 1.3}{Group Activity --- Answer Key}

\vspace{0.12in}

\begin{scenariobox}[The Water Balloon Launch]{plumlight}
\small
A water balloon is launched straight up from the ground. The function $h$
gives its height in feet $t$ seconds after launch. Every tier uses this same
table and pre-drawn graph --- no drawing required.

\smallskip
\begin{center}
\renewcommand{\arraystretch}{1.4}
\begin{tabular}{|c|c|c|c|c|c|}
\hline
$t$ (seconds) & 0 & 1 & 2 & 3 & 4 \\
\hline
$h(t)$ (feet) & 0 & 48 & 64 & 48 & 0 \\
\hline
\end{tabular}
\end{center}

\begin{center}
\begin{tikzpicture}[x=1.1cm, y=0.55cm]
  \draw[linegray, very thin, xstep=1, ystep=1] (0,0) grid (4,4);
  \draw[->, thick] (0,0) -- (4.5,0) node[below right] {\small $t$ (s)};
  \draw[->, thick] (0,0) -- (0,4.6) node[above] {\small $h(t)$ (ft)};
  \foreach \x in {1,2,3,4} \node[below] at (\x,0) {\small \x};
  \node[left] at (0,1) {\small 16};
  \node[left] at (0,2) {\small 32};
  \node[left] at (0,3) {\small 48};
  \node[left] at (0,4) {\small 64};
  \draw[very thick, plum] plot [smooth] coordinates
    {(0,0) (1,3) (2,4) (3,3) (4,0)};
  \fill[plum] (0,0) circle (2.5pt);
  \fill[plum] (1,3) circle (2.5pt);
  \fill[plum] (2,4) circle (2.5pt);
  \fill[plum] (3,3) circle (2.5pt);
  \fill[plum] (4,0) circle (2.5pt);
\end{tikzpicture}
\end{center}
\end{scenariobox}

\vspace{0.1in}

\boxguard
\begin{tcolorbox}[colback=white, colframe=black!40,
    title=\textbf{Tier R --- Remediate},
    fonttitle=\bfseries, arc=2mm, left=3mm, right=3mm, top=2mm, bottom=2mm]

\begin{enumerate}[itemsep=8pt]
  \item Fill in the blanks to compute the AROC over $[0, 1]$:

        \vspace{0.05in}
        AROC $= \dfrac{h(1) - h(0)}{1 - 0} = \dfrac{{\color{keyred}\mathbf{48}} - {\color{keyred}\mathbf{0}}}{1 - 0}
        = $ \ans{48} feet per second

  \item Now the AROC over $[1, 2]$:

        \vspace{0.05in}
        AROC $= \dfrac{h(2) - h(1)}{2 - 1} = $ \ans{16} feet per second

  \item During which second did the balloon climb \textbf{faster} --- the
        first ($0$ to $1$) or the second ($1$ to $2$)? \quad \ans{the first}

  \item Between $t = 2$ and $t = 3$, is the balloon rising or falling? \quad
        \textbf{Rising} / \ans{Falling}
\end{enumerate}
\end{tcolorbox}

\vspace{0.1in}

\boxguard
\begin{tcolorbox}[colback=white, colframe=black!40,
    title=\textbf{Tier A --- Approaching Proficiency},
    fonttitle=\bfseries, arc=2mm, left=3mm, right=3mm, top=2mm, bottom=2mm]

\begin{enumerate}[itemsep=8pt]
  \item Complete the AROC table.

        \smallskip
        \renewcommand{\arraystretch}{1.5}
        \begin{tabular}{|c|c|c|c|c|}
        \hline
        \textbf{Interval} & $[0,1]$ & $[1,2]$ & $[2,3]$ & $[3,4]$ \\
        \hline
        \textbf{AROC (ft/s)} & \ans{48} & \ans{16} & \ans{$-16$} & \ans{$-48$} \\
        \hline
        \end{tabular}

  \item Write one sentence interpreting the AROC over $[2, 3]$ --- include
        its sign and its units.

        \vspace{0.05in}
        \ansline{On average the balloon fell 16 feet each second between $t=2$ and $t=3$ ---}
        \ansline{the negative sign says height decreases as time increases.}

  \item Look at your table: each AROC changes by the same amount from one
        interval to the next. That step is \ans{$-32$} ft/s per second ---
        so $h$ is a \ans{quadratic} function, and its graph is concave
        \ans{down}.

  \item Compute the AROC over the whole flight, $[0, 4]$. \quad
        \ans{0} ft/s

        The balloon clearly moved --- why is this answer technically correct
        but misleading?

        \vspace{0.05in}
        \ansline{The net change is $0 - 0 = 0$: it ended at its starting height. The average}
        \ansline{erases the whole up-and-down story in between.}
\end{enumerate}
\end{tcolorbox}

\vspace{0.1in}

\boxguard
\begin{tcolorbox}[colback=white, colframe=black!40,
    title=\textbf{Tier E --- Extension},
    fonttitle=\bfseries, arc=2mm, left=3mm, right=3mm, top=2mm, bottom=2mm]

\begin{enumerate}[itemsep=8pt]
  \item How fast is the balloon moving \textit{at the instant} $t = 1$? An
        average needs two points --- so shrink the interval. Use these extra
        values:

        \smallskip
        \renewcommand{\arraystretch}{1.4}
        \begin{tabular}{|c|c|c|}
        \hline
        $t$ & 1.5 & 1.1 \\
        \hline
        $h(t)$ & 60 & 51.04 \\
        \hline
        \end{tabular}
        \smallskip

        \begin{enumerate}[label=(\alph*), itemsep=4pt]
          \item AROC over $[1, 2]$: \ans{16} ft/s
          \item AROC over $[1, 1.5]$: \ans{24} ft/s
          \item AROC over $[1, 1.1]$: \ans{30.4} ft/s
        \end{enumerate}

  \item The AROCs in question 1 are closing in on a value. Estimate the
        balloon's speed at the instant $t = 1$: about \ans{32} ft/s.
        Why does a smaller interval give a better estimate?

        \vspace{0.05in}
        \ansline{A smaller interval leaves less room for the speed to change inside it, so the}
        \ansline{average over the interval stays closer to the speed at the instant itself.}

  \item At what instant is the balloon's rate of change equal to zero, and
        what is happening physically at that moment?

        \vspace{0.05in}
        \ansline{At $t = 2$ seconds --- the top of the flight. The balloon has stopped rising}
        \ansline{and is about to fall, so for an instant its height is not changing.}
\end{enumerate}
\end{tcolorbox}

\end{document}
